% Motor: XeLaTeX, fijado en .latexmkrc y en la receta del editor.
\documentclass[aspectratio=169,10pt,c]{beamer}

% !TEX root = slides.tex
% Preambulo comun: estilo UTP, colores, fuentes, plantillas de beamer y
% comandos de apoyo didactico. Se carga desde slides.tex tras \documentclass.

% La opción "c" de beamer reparte el espacio libre 1:1.5 (arriba:abajo), por lo
% que el contenido queda por encima del centro. Igualamos ambos rellenos.
\makeatletter
\define@key{beamerframe}{c}[true]{%
  \beamer@frametopskip=0pt plus 1fill\relax%
  \beamer@framebottomskip=0pt plus 1fill\relax%
  \beamer@frametopskipautobreak=0pt plus .5\paperheight\relax%
  \beamer@framebottomskipautobreak=0pt plus .5\paperheight\relax%
}
\makeatother

\usepackage{graphicx}
\usepackage{tikz}
\usepackage{xcolor}
\usepackage{iftex}

\ifPDFTeX
  \usepackage[utf8]{inputenc}
  \usepackage[T1]{fontenc}
  \usepackage{lmodern}          % tipografia escalable para el respaldo con pdflatex
  \renewcommand{\familydefault}{\sfdefault}
\else
  \usepackage{fontspec}
  \IfFontExistsTF{Arial}
    {\setsansfont{Arial}[BoldFont={Arial Bold}]}
    {\setsansfont{Helvetica}[BoldFont={Helvetica Bold}]}
\fi

\usepackage[spanish,es-nodecimaldot]{babel}
\usepackage{microtype}
\usepackage{booktabs}
\usepackage{amssymb}
\usepackage{array}
\usetikzlibrary{arrows.meta,positioning,fit,calc,shapes.geometric}

\graphicspath{{assets/}}

% ---------- Identidad visual UTP ----------
\definecolor{utpRojo}{HTML}{C42B35}
\definecolor{utpNegro}{HTML}{111111}
\definecolor{utpGris}{HTML}{5F5F5F}
\definecolor{utpGrisMedio}{HTML}{9A9A9A}
\definecolor{utpGrisClaro}{HTML}{F1F1F1}
\definecolor{utpNaranja}{HTML}{C1670B}

\usetheme{default}
\setbeamercolor{normal text}{fg=utpNegro,bg=white}
\setbeamercolor{structure}{fg=utpRojo}
\setbeamercolor{frametitle}{fg=utpNegro,bg=white}
\setbeamercolor{title}{fg=utpNegro}
\setbeamercolor{subtitle}{fg=utpNegro}
\setbeamercolor{author}{fg=utpNegro}
\setbeamercolor{date}{fg=utpNegro}
\setbeamercolor{block title}{fg=white,bg=utpNegro}
\setbeamercolor{block body}{fg=utpNegro,bg=utpGrisClaro}
\setbeamercolor{alertblock title}{fg=white,bg=utpRojo}
\setbeamercolor{alertblock body}{fg=utpNegro,bg=utpRojo!7}
\setbeamercolor{exampleblock title}{fg=white,bg=utpGris}
\setbeamercolor{exampleblock body}{fg=utpNegro,bg=utpGris!8}

\setbeamerfont{title}{series=\bfseries,size=\fontsize{19.5}{22.5}\selectfont}
\setbeamerfont{subtitle}{size=\fontsize{11.5}{13.5}\selectfont}
\setbeamerfont{date}{series=\bfseries,size=\fontsize{11.5}{13.5}\selectfont}
\setbeamerfont{author}{size=\fontsize{10.5}{12.5}\selectfont}
\setbeamerfont{frametitle}{series=\bfseries,size=\large}
\setbeamerfont{block title}{series=\bfseries,size=\small}
\setbeamerfont{block body}{size=\normalsize}

\setbeamertemplate{navigation symbols}{}
\setbeamertemplate{headline}{}
\setbeamertemplate{blocks}[default]
\setbeamertemplate{itemize item}{\color{utpRojo}\small\raise.15ex\hbox{$\blacktriangleright$}}
\setbeamertemplate{itemize subitem}{\color{utpGris}\scriptsize$\bullet$}
\setbeamertemplate{enumerate item}{\color{utpRojo}\insertenumlabel.}
\setbeamertemplate{background canvas}{\color{white}\rule{\paperwidth}{\paperheight}}

% Título de diapositiva: negrita negra a la izquierda, logo UTP a la derecha
\setbeamertemplate{frametitle}{%
  \begin{tikzpicture}[remember picture,overlay]
    \node[anchor=north east,inner sep=0pt]
      at ([xshift=-7mm,yshift=-5.5mm]current page.north east)
      {\includegraphics[width=17mm]{utp-logo.png}};
  \end{tikzpicture}%
  \vspace{1.2mm}%
  \begin{minipage}[t]{.80\textwidth}\raggedright\insertframetitle\end{minipage}\par
  \vspace{1.1mm}{\color{utpRojo}\rule{\textwidth}{1.1pt}}
}

% Solo numeración discreta; la plantilla UTP no lleva pie de página
\setbeamertemplate{footline}{%
  \hfill{\color{utpGrisMedio}\scriptsize\insertframenumber}\hspace{6mm}\vspace{2.5mm}%
}

% Portada UTP
\setbeamertemplate{title page}{%
  \begin{tikzpicture}[remember picture,overlay]
    \node[anchor=center,inner sep=0pt] at (current page.center) {%
      \includegraphics[width=\paperwidth,height=\paperheight]{utp-cover-bg.jpeg}%
    };
    \node[anchor=north,align=center,text=utpNegro,inner sep=0pt,
      text width=145mm,font=\usebeamerfont{title}]
      at ([yshift=-25.2mm]current page.north) {\inserttitle};
    \node[anchor=north,align=center,text=utpNegro,inner sep=0pt,
      text width=142mm,font=\usebeamerfont{subtitle}]
      at ([yshift=-38.2mm]current page.north) {\insertsubtitle};
    \node[anchor=north,align=center,text=utpNegro,inner sep=0pt,
      font=\usebeamerfont{date}]
      at ([yshift=-50.4mm]current page.north) {\insertdate};
    \node[anchor=north,align=center,text=utpNegro,inner sep=0pt,
      font=\usebeamerfont{author}]
      at ([yshift=-58.4mm]current page.north) {\insertauthor};
  \end{tikzpicture}%
}

% ---------- Comandos de apoyo didáctico ----------
\newcommand{\concepto}[2]{%
  \begin{block}{#1}\large #2\end{block}}
\newcommand{\pregunta}[1]{%
  \begin{alertblock}{Decide y explica}#1\end{alertblock}}
\newcommand{\analogia}[1]{%
  \begin{exampleblock}{En palabras simples}\small #1\end{exampleblock}}
\newcommand{\enlarealidad}[1]{%
  \begin{block}{Cómo se ve en la realidad}\small #1\end{block}}
\newcommand{\ojo}[1]{%
  \begin{alertblock}{Ojo con esto}\small #1\end{alertblock}}
\newcommand{\apoyo}{{\normalfont\footnotesize\color{utpGris}~· material de apoyo}}
\newcommand{\marca}[1]{\tikz[baseline=(m.base)]{%
  \node[rounded corners=1.5pt,fill=utpGrisClaro,draw=utpGrisMedio,inner xsep=3pt,inner ysep=1.5pt,
        font=\footnotesize] (m) {#1};}}

\tikzset{
  caja/.style={draw=utpNegro, very thick, rounded corners=1.5pt, fill=white,
    align=center, minimum height=0.72cm, inner sep=5pt, font=\small},
  cajaazul/.style={caja, draw=utpGris, fill=white},
  cajaverde/.style={caja, draw=utpNegro, fill=utpGrisClaro},
  cajanaranja/.style={caja, draw=utpRojo!75, fill=utpRojo!7},
  flecha/.style={-{Stealth[length=2.3mm]}, very thick, draw=utpNegro},
  flechavuelta/.style={{Stealth[length=2.3mm]}-{Stealth[length=2.3mm]}, very thick, draw=utpNegro}
}

% ---------- Notación pata de gallo (crow's foot) ----------
% La rotación indica hacia dónde está la entidad: 0 derecha, 180 izquierda,
% 90 arriba, -90 abajo. La pata se abre hacia la entidad del lado "muchos".
\tikzset{
  relacion/.style={thick, draw=utpNegro},
  pics/uno/.style={code={%
    \draw[thick,utpNegro,line cap=round] (0,-1.5mm) -- (0,1.5mm);}},
  pics/muchos/.style={code={%
    \draw[thick,utpNegro,line cap=round,line join=round]
      (2.6mm,1.5mm) -- (0,0) -- (2.6mm,-1.5mm);
    \draw[thick,utpNegro,line cap=round] (0,0) -- (2.6mm,0);}},
}
\newcommand{\simbuno}{\tikz[baseline=-0.6ex]{\draw[relacion](0,0)--(7mm,0);\pic at (4.4mm,0){uno};}}
\newcommand{\simbmuchos}{\tikz[baseline=-0.6ex]{\draw[relacion](0,0)--(7mm,0);\pic at (4.4mm,0){muchos};}}

\usepackage{colortbl}

\title{Combinar consultas, subconsultas y CTE}
\subtitle{\texttt{UNION}, \texttt{INTERSECT}, \texttt{EXCEPT}, \texttt{EXISTS} y consultas con nombre}
\date{Semana 12 - Sesión 24}
\author{Rel Guzman}

% Nombre de tabla, siempre en naranja
\newcommand{\tbl}[1]{\textcolor{utpNaranja}{\texttt{#1}}}
% Palabra reservada de SQL dentro de un bloque de código
\newcommand{\kw}[1]{\textcolor{utpRojo}{\textbf{#1}}}
% Texto de apoyo en gris dentro de una tabla
\newcommand{\gris}[1]{\textcolor{utpGris}{#1}}
% Insignia de práctica
\newcommand{\ins}[1]{\tikz[baseline=-0.5ex]{\node[rounded corners=2pt,fill=utpRojo,text=white,
  font=\bfseries\scriptsize,inner xsep=5pt,inner ysep=2.5pt] {#1};}}
% Etiqueta oscura: nombre de un operador o de un gestor
\newcommand{\etiqueta}[1]{\tikz[baseline=-0.5ex]{\node[rounded corners=2pt,fill=utpNegro,text=white,
  font=\bfseries\scriptsize,inner xsep=6pt,inner ysep=3pt] {#1};}}
% Número de paso
\newcommand{\npieza}[1]{\tikz[baseline=-0.6ex]{\node[circle,fill=utpRojo,text=white,
  font=\bfseries\scriptsize,inner sep=1.6pt,minimum size=4.2mm] {#1};}}
% Celda resaltada en rojo claro
\newcommand{\mal}[1]{{\setlength{\fboxsep}{1.2pt}\colorbox{utpRojo!18}{#1}}}
% La subconsulta, resaltada en amarillo dentro de un bloque de código
\newcommand{\dentro}[1]{{\setlength{\fboxsep}{1.5pt}\colorbox{yellow!45}{#1}}}
% Fila de una tabla que cumple la condición
\newcommand{\fila}{\rowcolor{yellow!45}}
% Marcas de sí y no al lado de una fila
\newcommand{\si}{$\checkmark$}
\newcommand{\no}{\textcolor{utpRojo}{$\times$}}
% Bloque de código sobre fondo gris
\newsavebox{\sqlbuf}
\newenvironment{sqlcode}[1][.86\textwidth]
  {\begin{center}\begin{lrbox}{\sqlbuf}\begin{minipage}{#1}\ttfamily\footnotesize\vspace{6pt}\raggedright}
  {\vspace{6pt}\end{minipage}\end{lrbox}\colorbox{utpGrisClaro}{\usebox{\sqlbuf}}\end{center}}
% Sangría dentro de un bloque de código
\newcommand{\ind}{\hspace*{1.6em}}
% Anatomía de una instrucción: cada pieza del código con su pregunta al lado
\tikzset{
  pieza/.style={fill=utpGrisClaro,font=\ttfamily\footnotesize,anchor=west,inner sep=5pt,
    minimum height=0.62cm},
  para/.style={font=\small,anchor=west,text=utpNegro}
}
\newcommand{\piezax}{7.6}
\newcommand{\pieza}[3]{%
  \node[pieza] at (0,#1) {#2};
  \node[para] at (\piezax,#1) {\textcolor{utpRojo}{$\blacktriangleleft$}\ #3};}
% Diapositiva de sección: una palabra bien grande (\seccion[0]{...} la deja sin línea roja)
\newcommand{\seccion}[2][1]{%
\begin{frame}[plain]
  \begin{tikzpicture}[remember picture,overlay]
    \node[anchor=north east,inner sep=0pt]
      at ([xshift=-10mm,yshift=-8mm]current page.north east)
      {\includegraphics[width=43mm]{utp-logo.png}};
    \node[anchor=center,inner sep=0pt] at (current page.center)
      {\resizebox{100mm}{!}{\bfseries\color{utpNegro}#2}};
    \ifnum#1=1
    \node[anchor=north,inner sep=0pt] at ([yshift=-11mm]current page.center)
      {\color{utpRojo}\rule{100mm}{1.6pt}};
    \fi
  \end{tikzpicture}
\end{frame}}
% Cabecera de columna dentro de una tabla en monoespaciado
\newcommand{\cab}[1]{\normalfont\bfseries #1}
% Resultado de una consulta: \res{columnas}{cabecera}{filas}
\newcommand{\res}[3]{%
  \centering{\scriptsize\color{utpGris}resultado}\\[3pt]
  \footnotesize\ttfamily
  \begin{tabular}{#1}
  \toprule
  #2 \\
  \midrule
  #3
  \bottomrule
  \end{tabular}}
% Lo mismo, cuando debajo sigue más contenido del frame
\newcommand{\resg}[3]{{\res{#1}{#2}{#3}\par}}
% Tabla pequeña con su nombre arriba: \mini{nombre}{columnas}{cabecera}{filas}
\newcommand{\mini}[4]{%
  \begin{tabular}{@{}c@{}}
  \tbl{#1}\\[3pt]
  {\scriptsize\ttfamily
  \begin{tabular}{#2}
  \toprule
  #3 \\
  \midrule
  #4
  \bottomrule
  \end{tabular}}
  \end{tabular}}
% Forma de lo que devuelve una consulta: celdas con borde
\newcommand{\forma}[2]{{\ttfamily\footnotesize\renewcommand{\arraystretch}{1.25}%
  \begin{tabular}{#1}\hline #2\end{tabular}}}
% Tabla cita reducida a las columnas de los ejemplos de empleados
\newcommand{\tablacita}{%
  \mini{cita}{rr}{\cab{id} & \cab{empleado\_id}}{%
    1 & 1 \\ 2 & 2 \\ 3 & 1 \\ 4 & 1 \\ 5 & 2 \\ 6 & 1 \\ 7 & 2 \\}}
% Dos listas superpuestas: \venn{izquierda}{centro}{derecha}, con el color de cada zona
\colorlet{on}{yellow!45}
\colorlet{off}{white}
\newcommand{\vennbase}[4]{%
  \begin{tikzpicture}
    \fill[#1] (-0.75,0) circle (1.4);
    \fill[#3] (0.75,0) circle (1.4);
    \begin{scope}
      \clip (-0.75,0) circle (1.4);
      \fill[#2] (0.75,0) circle (1.4);
    \end{scope}
    \draw[very thick,utpNegro] (-0.75,0) circle (1.4);
    \draw[very thick,utpNegro] (0.75,0) circle (1.4);
    #4
  \end{tikzpicture}}
\newcommand{\venn}[3]{\vennbase{#1}{#2}{#3}{%
    \node[font=\scriptsize\bfseries,anchor=south] at (-1.2,1.42) {cliente};
    \node[font=\scriptsize\bfseries,anchor=south] at (1.2,1.42) {empleado};
    \node[font=\footnotesize,align=center] at (-1.42,0) {Lucía\\[1pt]Pedro\\[1pt]Rosa};
    \node[font=\footnotesize,align=center] at (0,0) {Bruno};
    \node[font=\footnotesize,align=center] at (1.42,0) {Ana\\[1pt]Carla};}}
% Versión pequeña, sin nombres
\newcommand{\vennchico}[3]{\scalebox{.62}{\vennbase{#1}{#2}{#3}{}}}
% Estilos de los diagramas de esta sesión
\tikzset{
  codigo/.style={font=\ttfamily\footnotesize,inner sep=3pt},
  sub/.style={codigo,draw=utpRojo,very thick,fill=yellow!45,rounded corners=2pt}
}
\begin{document}
{\setbeamertemplate{footline}{}\maketitle}
\begin{frame}[c]
\begin{block}{Logro de la sesión}
Al finalizar, podrás \textbf{combinar} los resultados de dos consultas, escribir \textbf{subconsultas} con \texttt{IN} y \texttt{EXISTS}, y ordenar una consulta larga por pasos con \textbf{CTE}, en MySQL y SQL Server.
\end{block}
\end{frame}
% =======================================================================
% Revisión de la sesión anterior
% ---------------------------------------------------------------------
\begin{frame}{Revisión de la sesión anterior: conectarse a un gestor}
\begin{columns}[c,onlytextwidth]
\column{.5\textwidth}
\small Para conectarte usas los \textbf{datos de conexión}:
\vspace{0.1cm}
\begin{center}
\footnotesize
\renewcommand{\arraystretch}{1.3}
\begin{tabular}{@{}>{\bfseries}ll@{}}
\toprule
Dato & \bfseries En clase \\
\midrule
Host & \texttt{localhost} \\
Puerto & \texttt{3306} \\
Usuario & \texttt{root} \\
Contraseña & la de la instalación \\
\bottomrule
\end{tabular}
\end{center}
\vspace{0.1cm}
\small También puedes iniciar la conexión hacia una \textbf{base de datos}, como \texttt{peluqueria}.
\column{.45\textwidth}
\small Si no conecta, revisa en este orden:
\vspace{0.2cm}
\begin{enumerate}\setlength\itemsep{7pt}
  \item ¿El \textbf{servicio} está encendido?
  \item ¿Los \textbf{datos de conexión} son correctos?
\end{enumerate}
\end{columns}
\vspace{0.4cm}
\concepto{Hoy}{\small Ya sabes conectarte. Hoy creamos \texttt{peluqueria} en los dos gestores y le hacemos preguntas más difíciles.}
\end{frame}
% ---------------------------------------------------------------------
\begin{frame}{\ins{Práctica 1} Crear la base en los dos gestores}
\small Los scripts del modelo físico están en \texttt{ejemplos/practicas/}.
\vspace{0.1cm}
\begin{enumerate}\setlength\itemsep{5pt}
  \item \textbf{MySQL:} en Workbench, abre y ejecuta \texttt{practica\_1\_modelo\_fisico\_mysql.sql}.
  \item \textbf{SQL Server:} en SSMS, abre y ejecuta \texttt{practica\_1\_modelo\_fisico\_sqlserver.sql}.
\end{enumerate}
\small Comprueba en cada gestor que la base se creó:
\begin{sqlcode}[.5\textwidth]
\kw{USE} peluqueria;\\
\kw{SELECT} nombre \kw{FROM} empleado;
\end{sqlcode}
\vspace{0.1cm}
\ojo{Para repetir en MySQL los ejemplos de clase, ejecuta también \texttt{ejemplos/consultas\_peluqueria/1\_datos\_de\_ejemplo.sql}. Este script \textbf{borra} los clientes, mascotas, citas y pagos, y los vuelve a cargar.}
\end{frame}
% =======================================================================
\seccion{Combinar SELECT}
\begin{frame}{Una pregunta real: ¿quién tiene descuento?}
\begin{center}
\begin{tikzpicture}
  \node[caja,cajanaranja,font=\large,inner sep=9pt] {¿Qué empleados también son clientes?};
\end{tikzpicture}
\end{center}
\vspace{0.1cm}
\small Empecemos con lo que ya sabes: una consulta por tabla.
\vspace{0.1cm}
\begin{columns}[T,onlytextwidth]
\column{.49\textwidth}
\centering\etiqueta{Clientes}
\begin{sqlcode}[\dimexpr\textwidth-2\fboxsep\relax]
\kw{SELECT} dni, nombre \kw{FROM} cliente;
\end{sqlcode}
\vspace{-0.1cm}
\res{ll}{\cab{dni} & \cab{nombre}}{45123678 & Lucía Ramos \\ 46234789 & Pedro Salas \\ 47345890 & Rosa Medina \\ 48456901 & Bruno Díaz \\}
\column{.49\textwidth}
\centering\etiqueta{Empleados}
\begin{sqlcode}[\dimexpr\textwidth-2\fboxsep\relax]
\kw{SELECT} dni, nombre \kw{FROM} empleado;
\end{sqlcode}
\vspace{-0.1cm}
\res{ll}{\cab{dni} & \cab{nombre}}{43567012 & Ana Torres \\ 48456901 & Bruno Díaz \\ 44678123 & Carla Núñez \\}
\end{columns}
\end{frame}
% ---------------------------------------------------------------------
\begin{frame}{Dos listas, tres preguntas}
\small Con esas dos listas puedes responder tres preguntas distintas, y cada una lleva a una decisión.
\vspace{0.2cm}
\begin{columns}[T,onlytextwidth]
\column{.31\textwidth}
\centering
\vennchico{on}{on}{on}\\[6pt]
\etiqueta{UNION}\\[6pt]
\small ¿Quién está en \textbf{alguna} de las dos listas?\\[6pt]
\footnotesize\gris{Enviarles el saludo de fin de año.}
\column{.31\textwidth}
\centering
\vennchico{off}{on}{off}\\[6pt]
\etiqueta{INTERSECT}\\[6pt]
\small ¿Quién está en \textbf{las dos}?\\[6pt]
\footnotesize\gris{Darles el descuento de personal.}
\column{.31\textwidth}
\centering
\vennchico{on}{off}{off}\\[6pt]
\etiqueta{EXCEPT}\\[6pt]
\small ¿Quién es cliente y \textbf{no} empleado?\\[6pt]
\footnotesize\gris{Enviarles la promoción de noviembre.}
\end{columns}
\vspace{0.3cm}
\concepto{Combinar SELECT}{\small Se escriben dos \texttt{SELECT} y, entre ellos, un operador. El resultado es \textbf{una sola lista}. \texttt{JOIN} agrega \textbf{columnas}; estos operadores trabajan con \textbf{filas}.}
\end{frame}
% ---------------------------------------------------------------------
\begin{frame}{\texttt{UNION}: todo lo de las dos listas}
\begin{columns}[c,onlytextwidth]
\column{.42\textwidth}
\centering
\venn{on}{on}{on}
\column{.55\textwidth}
\begin{sqlcode}[\dimexpr\textwidth-2\fboxsep\relax]
\kw{SELECT} dni, nombre \kw{FROM} cliente\\
\dentro{\kw{UNION}}\\
\kw{SELECT} dni, nombre \kw{FROM} empleado;
\end{sqlcode}
\vspace{-0.1cm}
\res{ll}{\cab{dni} & \cab{nombre}}{45123678 & Lucía Ramos \\ 46234789 & Pedro Salas \\ 47345890 & Rosa Medina \\ 48456901 & Bruno Díaz \\ 43567012 & Ana Torres \\ 44678123 & Carla Núñez \\}
\end{columns}
\vspace{0.3cm}
\analogia{Bruno Díaz está en las dos tablas con el mismo DNI y sale \textbf{una sola vez}: \texttt{UNION} quita los repetidos. Con \texttt{UNION ALL}, que los deja, saldría dos veces.}
\end{frame}
% ---------------------------------------------------------------------
\begin{frame}{\texttt{INTERSECT}: solo lo que está en las dos}
\begin{columns}[c,onlytextwidth]
\column{.42\textwidth}
\centering
\venn{off}{on}{off}
\column{.55\textwidth}
\begin{sqlcode}[\dimexpr\textwidth-2\fboxsep\relax]
\kw{SELECT} dni, nombre \kw{FROM} cliente\\
\dentro{\kw{INTERSECT}}\\
\kw{SELECT} dni, nombre \kw{FROM} empleado;
\end{sqlcode}
\vspace{-0.1cm}
\res{ll}{\cab{dni} & \cab{nombre}}{48456901 & Bruno Díaz \\}
\end{columns}
\vspace{0.3cm}
\analogia{Bruno Díaz está en \textbf{las dos listas} con el mismo DNI: le corresponde el descuento de personal. El DNI evita confundir a dos personas que se llaman igual.}
\end{frame}
% ---------------------------------------------------------------------
\begin{frame}{\texttt{EXCEPT}: lo que está en la primera y no en la segunda}
\begin{columns}[c,onlytextwidth]
\column{.42\textwidth}
\centering
\venn{on}{off}{off}
\column{.55\textwidth}
\begin{sqlcode}[\dimexpr\textwidth-2\fboxsep\relax]
\kw{SELECT} dni, nombre \kw{FROM} cliente\\
\dentro{\kw{EXCEPT}}\\
\kw{SELECT} dni, nombre \kw{FROM} empleado;
\end{sqlcode}
\vspace{-0.1cm}
\res{ll}{\cab{dni} & \cab{nombre}}{45123678 & Lucía Ramos \\ 46234789 & Pedro Salas \\ 47345890 & Rosa Medina \\}
\end{columns}
\vspace{0.3cm}
\analogia{Lucía, Pedro y Rosa son clientes y \textbf{no trabajan aquí}: reciben la promoción. \textbf{El orden importa}: al revés, empleado \texttt{EXCEPT} cliente da \textbf{Ana} y \textbf{Carla}.}
\end{frame}
% ---------------------------------------------------------------------
\begin{frame}{Reglas para combinar dos \texttt{SELECT}}
\begin{itemize}\setlength\itemsep{5pt}
  \item[\npieza{1}] La misma \textbf{cantidad de columnas} en los dos \texttt{SELECT}.
  \item[\npieza{2}] Las columnas en el mismo \textbf{orden} y con tipos compatibles.
  \item[\npieza{3}] Los \textbf{nombres} de las columnas salen del primer \texttt{SELECT}.
  \item[\npieza{4}] Un solo \texttt{ORDER BY}, al \textbf{final} de todo.
\end{itemize}
\begin{sqlcode}[.8\textwidth]
\kw{SELECT} nombre, telefono \kw{FROM} cliente \kw{UNION SELECT} nombre \kw{FROM} empleado;\\[4pt]
\color{utpRojo}ERROR 1222 (21000): The used SELECT statements have a different number of columns
\end{sqlcode}
\vspace{-0.2cm}
\ojo{\texttt{INTERSECT} y \texttt{EXCEPT} funcionan desde \textbf{MySQL 8.0.31}. Revisa tu versión con \texttt{SELECT VERSION();}. En versiones anteriores dan error de sintaxis y se reemplazan con \texttt{IN} y \texttt{NOT IN}.}
\end{frame}
% ---------------------------------------------------------------------
\begin{frame}{\ins{Práctica 2} \texttt{EXCEPT} en los dos gestores}
\small Abre \texttt{ejemplos/practicas/practica\_2\_combinar.sql} y ejecútalo en los dos gestores.
\vspace{0.1cm}
\begin{enumerate}\setlength\itemsep{3pt}
  \item \textbf{MySQL:} en Workbench.
  \item \textbf{SQL Server:} en SSMS.
\end{enumerate}
\vspace{0.1cm}
\small ¿Qué servicios cuestan menos de S/ 60, sin contar los que duran una hora o más?
\begin{columns}[c,onlytextwidth]
\column{.57\textwidth}
\begin{sqlcode}[\dimexpr\textwidth-2\fboxsep\relax]
\kw{USE} peluqueria;\\
\kw{SELECT} nombre, precio \kw{FROM} servicio\\
\kw{WHERE} precio < 60\\
\dentro{\kw{EXCEPT}}\\
\kw{SELECT} nombre, precio \kw{FROM} servicio\\
\kw{WHERE} duracion >= 60\\
\kw{ORDER BY} precio;
\end{sqlcode}
\column{.39\textwidth}
\res{lr}{\cab{nombre} & \cab{precio}}{Corte de uñas & 15.00 \\ Baño e higiene & 35.00 \\}
\end{columns}
\vspace{0.1cm}
\concepto{Idea clave}{\small La lista sale igual en los dos gestores: \texttt{EXCEPT} se escribe de la misma forma.}
\end{frame}
% =======================================================================
\seccion{Subconsultas}
\begin{frame}{Una pregunta que necesita otra pregunta}
\begin{center}
\begin{tikzpicture}
  \node[caja,cajanaranja,font=\large,inner sep=9pt] {¿Qué servicios cuestan \textbf{más que el promedio}?};
\end{tikzpicture}
\end{center}
\vspace{0.1cm}
\begin{columns}[T,onlytextwidth]
\column{.49\textwidth}
\npieza{1}\ \small Primero, ¿cuál es el promedio?
\begin{sqlcode}[\dimexpr\textwidth-2\fboxsep\relax]
\kw{SELECT AVG}(precio) \kw{FROM} servicio;
\end{sqlcode}
\vspace{-0.1cm}
\res{r}{\cab{AVG(precio)}}{46.250000 \\}
\column{.49\textwidth}
\npieza{2}\ \small Luego, ¿qué servicios cuestan más?
\begin{sqlcode}[\dimexpr\textwidth-2\fboxsep\relax]
\kw{SELECT} nombre, precio \kw{FROM} servicio\\
\kw{WHERE} precio > \mal{46.25};
\end{sqlcode}
\vspace{-0.1cm}
\res{lr}{\cab{nombre} & \cab{precio}}{Corte de pelo & 55.00 \\ Baño completo con corte & 80.00 \\}
\end{columns}
\vspace{0.2cm}
\ojo{Funciona, pero copiaste \texttt{46.25} \textbf{a mano}. Si mañana cambia un precio, ese número queda viejo.}
\end{frame}
% ---------------------------------------------------------------------
\begin{frame}{En una sola: la subconsulta}
\small En lugar de escribir el número, escribe \textbf{la consulta que lo calcula}:
\vspace{0.2cm}
\begin{center}
\begin{tikzpicture}
  \node[codigo,anchor=west] (a) at (0,0) {\kw{WHERE} precio >};
  \node[codigo,anchor=west,draw=utpRojo,thick,fill=utpRojo!18] (n) at ([xshift=2pt]a.east) {46.25};
  \node[codigo,anchor=west] (b) at (0,-1.5) {\kw{WHERE} precio >};
  \node[sub,anchor=west] (s) at ([xshift=2pt]b.east) {(\kw{SELECT AVG}(precio) \kw{FROM} servicio)};
  \draw[flecha] (n.south)--node[right,font=\scriptsize]{se reemplaza por}(n.south|-s.north);
\end{tikzpicture}
\end{center}
\vspace{0.1cm}
\begin{sqlcode}[.62\textwidth]
\kw{SELECT} nombre, precio\\
\kw{FROM} servicio\\
\kw{WHERE} precio > \dentro{(\kw{SELECT AVG}(precio) \kw{FROM} servicio)};
\end{sqlcode}
\vspace{-0.1cm}
\concepto{Subconsulta}{\small Una consulta escrita \textbf{dentro de otra}, entre \textbf{paréntesis}. El gestor resuelve \textbf{primero} la de adentro y usa su resultado en la de afuera.}
\end{frame}
% ---------------------------------------------------------------------
\begin{frame}{¿Qué devuelve la de adentro?}
\small Lo primero es mirar la \textbf{forma} de su resultado. Eso decide cómo se usa.
\vspace{0.3cm}
\begin{columns}[T,onlytextwidth]
\column{.31\textwidth}
\centering
{\large\bfseries Un valor}\\[8pt]
\forma{|c|}{\cellcolor{yellow!45}46.250000 \\ \hline}\\[8pt]
\footnotesize una fila, una columna\\[10pt]
\etiqueta{= \ > \ < }\\[6pt]
\scriptsize se compara
\column{.31\textwidth}
\centering
{\large\bfseries Una lista}\\[8pt]
\forma{|c|}{\cellcolor{yellow!45}1 \\ \hline \cellcolor{yellow!45}2 \\ \hline}\\[8pt]
\footnotesize una columna, varias filas\\[10pt]
\etiqueta{IN \ NOT IN}\\[6pt]
\scriptsize se busca dentro
\column{.31\textwidth}
\centering
{\large\bfseries Una tabla}\\[8pt]
\forma{|c|c|}{\cellcolor{yellow!45}1 & \cellcolor{yellow!45}4 \\ \hline \cellcolor{yellow!45}2 & \cellcolor{yellow!45}3 \\ \hline}\\[8pt]
\footnotesize varias columnas y filas\\[10pt]
\etiqueta{FROM}\\[6pt]
\scriptsize se usa como tabla
\end{columns}
\vspace{0.2cm}
\ojo{Si la de adentro devuelve una lista y usas \texttt{=}, sale \texttt{ERROR 1242: Subquery returns more than 1 row}. El signo \texttt{=} espera un solo valor; para una lista, usa \texttt{IN}.}
\end{frame}
% ---------------------------------------------------------------------
\begin{frame}{En \texttt{WHERE}: un valor de otra tabla}
\small ¿Qué citas tiene Ana Torres? En \tbl{cita} no está su nombre, solo su \texttt{empleado\_id}. \textbf{No sabes} ese id: que lo busque la subconsulta.
\vspace{0.2cm}
\begin{center}
\begin{tikzpicture}
  \node (em) at (0,0) {\mini{empleado}{rllr}{\cab{id} & \cab{nombre} & \cab{dni} & \cab{orden}}{%
    \fila 1 & Ana Torres & 43567012 & 1 \\
    2 & Bruno Díaz & 48456901 & 2 \\
    3 & Carla Núñez & 44678123 & 3 \\}};
  \node (re) at (8,0) {%
    \begin{tabular}{@{}c@{}}
    {\scriptsize\color{utpGris}resultado}\\[3pt]
    {\footnotesize\ttfamily
    \begin{tabular}{rl}
    \toprule
    \cab{id} & \cab{fecha\_hora} \\
    \midrule
    1 & 2026-09-08 10:00:00 \\
    3 & 2026-09-22 09:00:00 \\
    4 & 2026-10-06 10:00:00 \\
    6 & 2026-10-20 09:00:00 \\
    \bottomrule
    \end{tabular}}
    \end{tabular}};
  \draw[flecha] (em.east)--node[above,font=\small\bfseries]{\npieza{1} id = 1}node[below,font=\scriptsize]{\npieza{2} busca en \texttt{cita}}(re.west);
\end{tikzpicture}
\end{center}
\vspace{0.1cm}
\begin{sqlcode}[.84\textwidth]
\kw{SELECT} id, fecha\_hora \kw{FROM} cita\\
\kw{WHERE} empleado\_id = \dentro{(\kw{SELECT} id \kw{FROM} empleado \kw{WHERE} nombre = 'Ana Torres')};
\end{sqlcode}
\end{frame}
% ---------------------------------------------------------------------
\begin{frame}{Con \texttt{IN} y \texttt{NOT IN}: ¿qué empleados tienen citas?}
\small \texttt{IN} ya lo conoces: \texttt{WHERE id IN (1, 2)}. Ahora la lista la escribe \textbf{una consulta}.
\begin{center}
\begin{tikzpicture}
  \node (ci) at (0,0) {\tablacita};
  \node (li) at (3.6,0) {\scalebox{.8}{\forma{|c|}{\cellcolor{yellow!45}1 \\ \hline \cellcolor{yellow!45}1 \\ \hline \cellcolor{yellow!45}1 \\ \hline \cellcolor{yellow!45}1 \\ \hline \cellcolor{yellow!45}2 \\ \hline \cellcolor{yellow!45}2 \\ \hline \cellcolor{yellow!45}2 \\ \hline}}};
  \node[font=\scriptsize\bfseries,above=2pt of li] {la lista};
  \node (em) at (7.6,0) {\mini{empleado}{rll}{\cab{id} & \cab{nombre} & }{%
    \fila 1 & Ana Torres & \si \\
    \fila 2 & Bruno Díaz & \si \\
    3 & Carla Núñez & \no \\}};
  \draw[flecha] (ci.east)--node[above,font=\scriptsize]{\npieza{1}}(li.west);
  \draw[flecha] (li.east)--node[above,font=\scriptsize]{\npieza{2}}node[below,font=\scriptsize]{¿está?}(em.west);
\end{tikzpicture}
\end{center}
\begin{columns}[T,onlytextwidth]
\column{.49\textwidth}
\centering
\begin{sqlcode}[\dimexpr\textwidth-2\fboxsep\relax]
\kw{SELECT} nombre \kw{FROM} empleado\\
\kw{WHERE} id \kw{IN}\\
\ind\dentro{(\kw{SELECT} empleado\_id \kw{FROM} cita)};
\end{sqlcode}
\vspace{-0.1cm}
\small \textbf{Ana Torres} y \textbf{Bruno Díaz}
\column{.49\textwidth}
\centering
\begin{sqlcode}[\dimexpr\textwidth-2\fboxsep\relax]
\kw{SELECT} nombre \kw{FROM} empleado\\
\kw{WHERE} id \kw{NOT IN}\\
\ind\dentro{(\kw{SELECT} empleado\_id \kw{FROM} cita)};
\end{sqlcode}
\vspace{-0.1cm}
\small \textbf{Carla Núñez}: aún no tiene citas
\end{columns}
\end{frame}
% ---------------------------------------------------------------------
\begin{frame}{En \texttt{FROM}: ¿cuántas citas tiene el que más tiene?}
\small \texttt{MAX(COUNT(*))} \textbf{no se puede} escribir. Primero cuentas; luego buscas el máximo.
\begin{center}
\begin{tikzpicture}
  \node (ci) at (0,0) {\tablacita};
  \node[draw=utpRojo,very thick,dashed,rounded corners=3pt,fill=yellow!20,inner sep=6pt] (co) at (4.6,0)
    {\mini{conteo}{rr}{\cab{empleado\_id} & \cab{total}}{1 & 4 \\ 2 & 3 \\}};
  \node[font=\scriptsize,text=utpRojo,below=2pt of co] {tabla temporal};
  \node[caja,cajaverde,font=\large\bfseries,minimum width=1.2cm] (r) at (8.6,0) {4};
  \draw[flecha] (ci.east)--node[above,font=\scriptsize]{\npieza{1} contar}(co.west);
  \draw[flecha] (co.east)--node[above,font=\scriptsize]{\npieza{2} \texttt{MAX}}(r.west);
\end{tikzpicture}
\end{center}
\vspace{0.1cm}
\begin{sqlcode}[.62\textwidth]
\kw{SELECT MAX}(total)\\
\kw{FROM} \dentro{(\kw{SELECT} empleado\_id, \kw{COUNT}(*) \kw{AS} total}\\
\ind\ind\dentro{\kw{FROM} cita \kw{GROUP BY} empleado\_id)} \kw{AS} conteo;
\end{sqlcode}
\vspace{-0.5cm}
\ojo{En \texttt{FROM}, la subconsulta \textbf{necesita un nombre} (\texttt{AS conteo}). Sin él: \texttt{ERROR 1248: Every derived table must have its own alias}.}
\end{frame}
% ---------------------------------------------------------------------
\begin{frame}{En \texttt{SELECT}: una columna más}
\small Si la de adentro devuelve \textbf{un valor}, puede ir como una columna. ¿Cuántas citas tiene \textbf{cada} empleado? Aquí la de adentro usa \texttt{e.id}, que viene de la consulta de afuera.
\begin{sqlcode}[.74\textwidth]
\kw{SELECT} e.nombre,\\
\ind\dentro{(\kw{SELECT COUNT}(*) \kw{FROM} cita c \kw{WHERE} c.empleado\_id = \mal{e.id})} \kw{AS} citas\\
\kw{FROM} empleado e;
\end{sqlcode}
\vspace{-0.1cm}
\begin{center}
\footnotesize
\renewcommand{\arraystretch}{1.2}
\begin{tabular}{@{}lll@{}}
\toprule
\bfseries Fila de afuera & \bfseries La de adentro cuenta\ldots & \bfseries citas \\
\midrule
Ana Torres, \texttt{id = 1} & las citas con \texttt{empleado\_id = 1} & \textbf{4} \\
Bruno Díaz, \texttt{id = 2} & las citas con \texttt{empleado\_id = 2} & \textbf{3} \\
Carla Núñez, \texttt{id = 3} & las citas con \texttt{empleado\_id = 3} & \textbf{0} \\
\bottomrule
\end{tabular}
\end{center}
\vspace{0.1cm}
\concepto{Subconsulta correlacionada}{\small Usa una columna de afuera. Por eso \textbf{no} se resuelve una sola vez: se ejecuta \textbf{una vez por cada fila} de afuera.}
\end{frame}
% ---------------------------------------------------------------------
\begin{frame}{\texttt{EXISTS}: ¿hay al menos una fila?}
\small No necesitas un dato. Solo quieres saber si \textbf{existe}: ¿este empleado tiene alguna cita?
\begin{sqlcode}[.76\textwidth]
\kw{SELECT} e.nombre\\
\kw{FROM} empleado e\\
\kw{WHERE EXISTS} \dentro{(\kw{SELECT} 1 \kw{FROM} cita c \kw{WHERE} c.empleado\_id = e.id)};
\end{sqlcode}
\vspace{-0.1cm}
\begin{center}
\footnotesize
\renewcommand{\arraystretch}{1.2}
\begin{tabular}{@{}llc@{}}
\toprule
\bfseries Fila de afuera & \bfseries ¿Tiene alguna cita? & \bfseries ¿Se queda? \\
\midrule
Ana Torres, \texttt{id = 1} & sí: encontró la cita 1 y dejó de buscar & \si \\
Bruno Díaz, \texttt{id = 2} & sí: encontró la cita 2 y dejó de buscar & \si \\
Carla Núñez, \texttt{id = 3} & no: revisó todas y no hay & \no \\
\bottomrule
\end{tabular}
\end{center}
\vspace{0.1cm}
\concepto{\texttt{EXISTS}}{\small Responde \textbf{sí o no}: sí, cuando la de adentro encuentra \textbf{al menos una fila}. No mira los valores, por eso se suele escribir \texttt{SELECT 1}.}
\end{frame}
% ---------------------------------------------------------------------
\begin{frame}{\texttt{NOT EXISTS}: que no haya ninguna}
\small Es la pregunta contraria. Sirve para encontrar lo que \textbf{quedó sin relación}.
\begin{columns}[T,onlytextwidth]
\column{.48\textwidth}
\centering\etiqueta{Empleados sin citas}
\begin{sqlcode}[\dimexpr\textwidth-2\fboxsep\relax]
\kw{SELECT} e.nombre\\
\kw{FROM} empleado e\\
\kw{WHERE NOT EXISTS} \dentro{(\kw{SELECT} 1}\\
\ind\dentro{\kw{FROM} cita c}\\
\ind\dentro{\kw{WHERE} c.empleado\_id = e.id)};
\end{sqlcode}
\vspace{-0.1cm}
\res{l}{\cab{nombre}}{Carla Núñez \\}
\column{.48\textwidth}
\centering\etiqueta{Citas atendidas sin cobrar}
\begin{sqlcode}[\dimexpr\textwidth-2\fboxsep\relax]
\kw{SELECT} c.id, c.fecha\_hora\\
\kw{FROM} cita c\\
\kw{WHERE} c.estado = 'completada'\\
\ind\kw{AND NOT EXISTS} \dentro{(\kw{SELECT} 1}\\
\ind\ind\dentro{\kw{FROM} pago p}\\
\ind\ind\dentro{\kw{WHERE} p.cita\_id = c.id)};
\end{sqlcode}
\vspace{-0.1cm}
\res{rl}{\cab{id} & \cab{fecha\_hora}}{6 & 2026-10-20 09:00:00 \\}
\end{columns}
\vspace{0.1cm}
\analogia{La segunda también sale con \texttt{EXCEPT}: citas completadas menos citas en \texttt{pago}. Con \texttt{NOT EXISTS} puedes mostrar más columnas, como la fecha.}
\end{frame}
% ---------------------------------------------------------------------
\begin{frame}{\texttt{EXISTS} vs \texttt{IN}}
\small Las dos consultas dan el mismo resultado: \textbf{Ana Torres} y \textbf{Bruno Díaz}.
\vspace{-0.1cm}
\begin{columns}[T,onlytextwidth]
\column{.45\textwidth}
\centering\etiqueta{IN}
\begin{sqlcode}[\dimexpr\textwidth-2\fboxsep\relax]
\kw{SELECT} nombre \kw{FROM} empleado\\
\kw{WHERE} id \kw{IN}\\
\ind\dentro{(\kw{SELECT} empleado\_id}\\
\ind\dentro{\kw{FROM} cita)};
\end{sqlcode}
\column{.52\textwidth}
\centering\etiqueta{EXISTS}
\begin{sqlcode}[\dimexpr\textwidth-2\fboxsep\relax]
\kw{SELECT} e.nombre \kw{FROM} empleado e\\
\kw{WHERE EXISTS}\\
\ind\dentro{(\kw{SELECT} 1 \kw{FROM} cita c}\\
\ind\dentro{\kw{WHERE} c.empleado\_id = e.id)};
\end{sqlcode}
\end{columns}
\vspace{-0.1cm}
\begin{center}
\footnotesize
\renewcommand{\arraystretch}{1.25}
\begin{tabular}{@{}>{\bfseries}l>{\raggedright}p{4.7cm}>{\raggedright\arraybackslash}p{5.1cm}@{}}
\toprule
 & \texttt{IN} & \texttt{EXISTS} \\
\midrule
Pregunta & ¿este valor está en la lista? & ¿hay al menos una fila? \\
La de adentro & se puede ejecutar sola & usa una columna de afuera \\
Cómo trabaja & arma la lista y busca en ella & prueba fila por fila \\
Conviene si & comparas una columna con una lista & necesitas varias condiciones \\
\bottomrule
\end{tabular}
\end{center}
\vspace{-0.1cm}
\analogia{Hoy los gestores suelen ejecutar las dos igual de rápido: elige la que se lea mejor. La diferencia real aparece con \texttt{NOT IN}.}
\end{frame}
% ---------------------------------------------------------------------
\begin{frame}{Ojo: \texttt{NOT IN} y los \texttt{NULL}}
\begin{columns}[T,onlytextwidth]
\column{.48\textwidth}
\begin{sqlcode}[\dimexpr\textwidth-2\fboxsep\relax]
\kw{SELECT} nombre \kw{FROM} empleado\\
\kw{WHERE} id \kw{NOT IN} (1, 2);
\end{sqlcode}
\vspace{-0.1cm}
\res{l}{\cab{nombre}}{Carla Núñez \\}
\column{.48\textwidth}
\begin{sqlcode}[\dimexpr\textwidth-2\fboxsep\relax]
\kw{SELECT} nombre \kw{FROM} empleado\\
\kw{WHERE} id \kw{NOT IN} (1, 2, \mal{NULL});
\end{sqlcode}
\vspace{-0.1cm}
\centering{\scriptsize\color{utpGris}resultado}\\[6pt]
\small\textbf{ninguna fila}
\end{columns}
\vspace{0.3cm}
\small ¿Por qué? Para Carla, el gestor compara su id con cada valor de la lista:
\vspace{0.1cm}
\begin{center}
\small
\texttt{3 <> 1} \si \qquad \texttt{3 <> 2} \si \qquad \texttt{3 <> NULL} \textbf{no se sabe}
\end{center}
\vspace{0.1cm}
\ojo{\texttt{NULL} es un valor desconocido. Si la lista tiene un \texttt{NULL}, el gestor no puede asegurar que el valor no esté, y descarta la fila. Pasa con columnas que admiten \texttt{NULL}, como \texttt{cliente.telefono}. \texttt{NOT EXISTS} no tiene este problema.}
\end{frame}
% ---------------------------------------------------------------------
\begin{frame}{\ins{Práctica 3} Subconsulta en los dos gestores}
\small Abre \texttt{ejemplos/practicas/practica\_3\_subconsulta.sql} y ejecútalo en los dos gestores.
\vspace{0.1cm}
\begin{enumerate}\setlength\itemsep{3pt}
  \item \textbf{MySQL:} en Workbench.
  \item \textbf{SQL Server:} en SSMS.
\end{enumerate}
\vspace{0.1cm}
\begin{columns}[c,onlytextwidth]
\column{.57\textwidth}
\begin{sqlcode}[\dimexpr\textwidth-2\fboxsep\relax]
\kw{USE} peluqueria;\\
\kw{SELECT} nombre, precio\\
\kw{FROM} servicio\\
\kw{WHERE} precio > \dentro{(\kw{SELECT AVG}(precio) \kw{FROM} servicio)};
\end{sqlcode}
\column{.39\textwidth}
\res{lr}{\cab{nombre} & \cab{precio}}{Corte de pelo & 55.00 \\ Baño completo con corte & 80.00 \\}
\end{columns}
\vspace{0.2cm}
\concepto{Idea clave}{\small La subconsulta también se escribe igual en los dos gestores.}
\end{frame}
% =======================================================================
\seccion{Expresiones de tabla comunes}
\begin{frame}{Cuando la consulta crece}
\small ¿Qué empleado tiene \textbf{más citas}? Con lo que sabemos, se puede escribir así:
\begin{sqlcode}[.96\textwidth]
\kw{SELECT} e.nombre, conteo.total\\
\kw{FROM} empleado e\\
\kw{JOIN} \dentro{(\kw{SELECT} empleado\_id, \kw{COUNT}(*) \kw{AS} total \kw{FROM} cita \kw{GROUP BY} empleado\_id)} \kw{AS} conteo\\
\ind\kw{ON} conteo.empleado\_id = e.id\\
\kw{WHERE} conteo.total = (\kw{SELECT MAX}(total)\\
\ind\kw{FROM} \dentro{(\kw{SELECT} empleado\_id, \kw{COUNT}(*) \kw{AS} total \kw{FROM} cita \kw{GROUP BY} empleado\_id)} \kw{AS} otro);
\end{sqlcode}
\vspace{0.1cm}
\small Funciona y devuelve \textbf{Ana Torres, 4}. Pero tiene dos problemas:
\vspace{0.1cm}
\begin{itemize}\setlength\itemsep{6pt}
  \item El mismo conteo está escrito \textbf{dos veces}.
  \item Se lee de \textbf{adentro hacia afuera}: hay que saltar por todo el código para entenderla.
\end{itemize}
\end{frame}
% ---------------------------------------------------------------------
\begin{frame}{CTE: una subconsulta con nombre}
\concepto{CTE (\textit{Common Table Expression})}{\small Una subconsulta a la que le pones \textbf{un nombre}, escrita \textbf{antes} de la consulta principal. En español: expresión de tabla común.}
\vspace{0.1cm}
\begin{center}
\renewcommand{\piezax}{7.2}
\begin{tikzpicture}
  \pieza{0}{\kw{WITH} conteo \kw{AS} (}{le pones un nombre}
  \pieza{-0.68}{\ind\kw{SELECT} empleado\_id, \kw{COUNT}(*) \kw{AS} total}{cuenta las citas\ldots}
  \pieza{-1.36}{\ind\kw{FROM} cita \kw{GROUP BY} empleado\_id}{\ldots{} de cada empleado}
  \pieza{-2.04}{)}{se cierra, sin punto y coma}
  \pieza{-2.72}{\kw{SELECT MAX}(total) \kw{FROM} conteo;}{la consulta la usa como tabla}
\end{tikzpicture}
\end{center}
\end{frame}
% ---------------------------------------------------------------------
\begin{frame}{Se escribe una vez y se usa dos}
\small Volvamos a la pregunta larga: ¿qué empleado tiene más citas?
\vspace{0.1cm}
\begin{sqlcode}[.72\textwidth]
\kw{WITH} conteo \kw{AS} \dentro{(\kw{SELECT} empleado\_id, \kw{COUNT}(*) \kw{AS} total}\\
\ind\ind\ind\dentro{\kw{FROM} cita \kw{GROUP BY} empleado\_id)}\\
\kw{SELECT} e.nombre, conteo.total\\
\kw{FROM} empleado e\\
\kw{JOIN} \mal{conteo} \kw{ON} conteo.empleado\_id = e.id\\
\kw{WHERE} conteo.total = (\kw{SELECT MAX}(total) \kw{FROM} \mal{conteo});
\end{sqlcode}
\vspace{-0.1cm}
\resg{lr}{\cab{nombre} & \cab{total}}{Ana Torres & 4 \\}
\vspace{0.3cm}
\concepto{Idea clave}{\small Si cambias cómo se cuenta, lo cambias en \textbf{un solo lugar}. Un CTE solo existe durante su consulta: \textbf{no} queda guardado en la base.}
\end{frame}
% ---------------------------------------------------------------------
\begin{frame}{Todo junto: la lista de cobros}
\small \texttt{WITH} se escribe \textbf{una vez} y los CTE se separan con \textbf{comas}. Cada uno puede usar los anteriores.
\begin{columns}[c,onlytextwidth]
\column{.6\textwidth}
\begin{sqlcode}[\dimexpr\textwidth-2\fboxsep\relax]
\kw{WITH} atendidas \kw{AS} (\kw{SELECT} id \kw{FROM} cita\\
\ind\ind\ind\ind\kw{WHERE} estado = 'completada'),\\
cobradas \kw{AS} (\kw{SELECT} cita\_id \kw{FROM} pago),\\
sin\_cobrar \kw{AS} (\kw{SELECT} id \kw{FROM} atendidas\\
\ind\ind\ind\ind\kw{EXCEPT}\\
\ind\ind\ind\ind\kw{SELECT} cita\_id \kw{FROM} cobradas)\\
\kw{SELECT} c.id, c.fecha\_hora,\\
\ind s.nombre \kw{AS} servicio, s.precio\\
\kw{FROM} sin\_cobrar sc\\
\kw{JOIN} cita c \kw{ON} c.id = sc.id\\
\kw{JOIN} servicio s \kw{ON} s.id = c.servicio\_id;
\end{sqlcode}
\column{.37\textwidth}
\begin{tikzpicture}[font=\scriptsize]
  \tikzset{e/.style={caja,font=\scriptsize,text width=3.4cm,minimum height=.55cm,inner sep=3pt}}
  \node[e] (a) at (0,0) {\npieza{1} citas atendidas};
  \node[e] (b) at (0,-.95) {\npieza{2} citas cobradas};
  \node[e] (c) at (0,-1.9) {\npieza{3} las que faltan cobrar};
  \node[cajanaranja,e] (d) at (0,-2.85) {\npieza{4} fecha, servicio y precio};
  \draw[flecha] (a)--(b); \draw[flecha] (b)--(c); \draw[flecha] (c)--(d);
\end{tikzpicture}
\end{columns}
\vspace{0.1cm}
\resg{rllr}{\cab{id} & \cab{fecha\_hora} & \cab{servicio} & \cab{precio}}{6 & 2026-10-20 09:00:00 & Corte de pelo & 55.00 \\}
\end{frame}
% =======================================================================
\begin{frame}{\ins{Práctica 4} CTE en los dos gestores}
\small Abre \texttt{ejemplos/practicas/practica\_4\_cte.sql} y ejecútalo en los dos gestores.
\vspace{0.1cm}
\begin{enumerate}\setlength\itemsep{3pt}
  \item \textbf{MySQL:} en Workbench.
  \item \textbf{SQL Server:} en SSMS.
\end{enumerate}
\vspace{0.1cm}
\begin{columns}[c,onlytextwidth]
\column{.57\textwidth}
\begin{sqlcode}[\dimexpr\textwidth-2\fboxsep\relax]
\kw{USE} peluqueria;\\
\kw{WITH} caros \kw{AS}\\
\ind\dentro{(\kw{SELECT} nombre, precio \kw{FROM} servicio}\\
\ind\dentro{\kw{WHERE} precio >}\\
\ind\ind\dentro{(\kw{SELECT AVG}(precio) \kw{FROM} servicio))}\\
\kw{SELECT COUNT}(*) \kw{AS} cantidad, \kw{SUM}(precio) \kw{AS} total\\
\kw{FROM} caros;
\end{sqlcode}
\column{.39\textwidth}
\res{rr}{\cab{cantidad} & \cab{total}}{2 & 135.00 \\}
\end{columns}
\vspace{0.2cm}
\concepto{Idea clave}{\small Es la consulta de la práctica 3, ahora con nombre. El CTE también se escribe igual en los dos gestores.}
\end{frame}
% ---------------------------------------------------------------------
% La misma escena en las dos láminas: la base a la izquierda y la app a la derecha.
\newcommand{\escenaapp}{%
  \tikzset{
    panel/.style={draw=utpNegro,very thick,fill=white,minimum height=3.1cm,anchor=north},
    barra/.style={text=white,minimum height=.58cm,anchor=north,
      font=\footnotesize\bfseries,inner sep=0pt}}
  % La consulta, arriba y al centro
  \node[sub,inner sep=6pt,font=\ttfamily\footnotesize] (sql) at (6.7,2.45)
    {\kw{SELECT} nombre, precio \kw{FROM} servicio;};
  % La base de datos
  \node[panel,minimum width=7cm] (bd) at (3.5,1.6) {};
  \node[barra,fill=utpGris,minimum width=7cm] (bdbarra) at (bd.north) {Base de datos: tabla \texttt{servicio}};
  \coordinate (ico) at ([xshift=.34cm]bdbarra.west);
  \draw[white,thick] ($(ico)+(-.15,.09)$) -- ++(0,-.18)
    arc[start angle=180,end angle=360,x radius=.15,y radius=.06] -- ++(0,.18);
  \draw[white,thick] ($(ico)+(0,.09)$) ellipse (.15 and .06);
  \node[anchor=north,inner sep=0pt] at ([yshift=-7pt]bdbarra.south) {%
    \scriptsize\ttfamily\renewcommand{\arraystretch}{1.42}\setlength{\tabcolsep}{3pt}%
    \begin{tabular}{r>{\columncolor{yellow!45}}lr>{\columncolor{yellow!45}}rr}
    \cab{id} & \cab{nombre} & \cab{duracion} & \cab{precio} & \cab{orden} \\ \hline
    1 & Baño e higiene & 45 & 35.00 & 1 \\
    2 & Corte de pelo & 60 & 55.00 & 2 \\
    3 & Corte de uñas & 15 & 15.00 & 3 \\
    4 & Baño completo con corte & 90 & 80.00 & 4 \\
    \end{tabular}};
  % El aplicativo
  \node[panel,minimum width=5.1cm] (app) at (10.85,1.6) {};
  \node[barra,fill=utpNegro,minimum width=5.1cm] (appbarra) at (app.north) {\hspace{.5cm}Peluquería: nuestros servicios};
  \foreach \i in {0,1,2}
    \fill[white] ([xshift={.2cm+\i*.17cm}]appbarra.west) circle (.05);
}
\newcommand{\vacio}[1]{\textcolor{gray!30}{\rule{#1}{1.5ex}}}
% ---------------------------------------------------------------------
\begin{frame}{Para esto sirven los \texttt{SELECT}: la app pide}
\begin{center}
\begin{tikzpicture}
  \escenaapp
  \node[anchor=north,inner sep=0pt] at ([yshift=-7pt]appbarra.south) {%
    \footnotesize\renewcommand{\arraystretch}{1.2}\setlength{\tabcolsep}{4pt}%
    \begin{tabular}{lr}
    \textbf{Servicio} & \textbf{Precio} \\ \hline
    \vacio{3.2cm} & \vacio{1.2cm} \\
    \vacio{3.2cm} & \vacio{1.2cm} \\
    \vacio{3.2cm} & \vacio{1.2cm} \\
    \vacio{3.2cm} & \vacio{1.2cm} \\
    \end{tabular}};
  % El recorrido
  \draw[flecha,draw=utpRojo] (12.5,1.6) |- (sql.east);
  \draw[flecha,draw=utpRojo] (sql.west) -| (1.2,1.6);
  \node[font=\scriptsize,text=utpRojo,anchor=south] at (11.3,2.47) {\npieza{1} la app pide};
\end{tikzpicture}
\end{center}
\vspace{0.1cm}
\concepto{La app pregunta}{\small Abres la pantalla de servicios y todavía está vacía: la app le envía este \texttt{SELECT} a la base. Solo pide las dos columnas que va a mostrar: \texttt{nombre} y \texttt{precio}.}
\end{frame}
% ---------------------------------------------------------------------
\begin{frame}{Para esto sirven los \texttt{SELECT}: la base responde}
\begin{center}
\begin{tikzpicture}
  \escenaapp
  \node[anchor=north,inner sep=0pt] at ([yshift=-7pt]appbarra.south) {%
    \footnotesize\renewcommand{\arraystretch}{1.2}\setlength{\tabcolsep}{4pt}%
    \begin{tabular}{>{\columncolor{yellow!45}}l>{\columncolor{yellow!45}}r}
    \textbf{Servicio} & \textbf{Precio} \\ \hline
    Baño e higiene & S/ 35.00 \\
    Corte de pelo & S/ 55.00 \\
    Corte de uñas & S/ 15.00 \\
    Baño completo con corte & S/ 80.00 \\
    \end{tabular}};
  % El recorrido
  \draw[flecha,draw=utpRojo] (7,-.05) -- (8.3,-.05);
  \node[font=\scriptsize,text=utpRojo,align=center,anchor=south] at (7.65,.03) {\npieza{2}\\[1pt]la base\\devuelve\\las filas};
\end{tikzpicture}
\end{center}
\vspace{0.1cm}
\concepto{Ese es el uso de los SELECT}{\small La tabla que ves en un aplicativo no se escribe a mano: sale de la base de datos con un \texttt{SELECT}. Si cambia un precio en la base, cambia en la pantalla.}
\end{frame}
\begin{frame}[plain]
  \begin{tikzpicture}[remember picture,overlay]
    \node[anchor=north west,font=\bfseries\fontsize{24}{28}\selectfont]
      at ([xshift=12mm,yshift=-9mm]current page.north west) {Preguntas};
    \node[anchor=north east,inner sep=0pt]
      at ([xshift=-10mm,yshift=-8mm]current page.north east)
      {\includegraphics[width=43mm]{utp-logo.png}};
    \node[anchor=center,inner sep=0pt]
      at ([yshift=-5mm]current page.center)
      {\includegraphics[width=78mm]{preguntas-ilustracion.png}};
  \end{tikzpicture}
\end{frame}
\begin{frame}{Cierre}
\begin{center}
\small
\renewcommand{\arraystretch}{1.15}
\begin{tabular}{>{\raggedright\bfseries}p{3.4cm}>{\raggedright\arraybackslash}p{8.8cm}}
\toprule
\texttt{UNION} & Junta las filas de dos \texttt{SELECT} sin repetir. \texttt{UNION ALL} deja los repetidos. \\
\texttt{INTERSECT} & Solo las filas que están en las dos listas. \\
\texttt{EXCEPT} & Las filas de la primera que no están en la segunda. El orden importa. \\
Subconsulta & Una consulta dentro de otra, entre paréntesis. Puede ir en \texttt{WHERE}, \texttt{FROM} o \texttt{SELECT}. \\
\texttt{IN} & Busca un valor en una lista. Cuidado con \texttt{NOT IN} y \texttt{NULL}. \\
\texttt{EXISTS} & Responde sí o no: ¿hay al menos una fila? \\
CTE & Una subconsulta con nombre, escrita antes con \texttt{WITH}. Solo vive durante su consulta. \\
En una app & Cada pantalla que muestra datos ejecuta un \texttt{SELECT}. \\
\bottomrule
\end{tabular}
\end{center}
\end{frame}

\begin{frame}{Referencias}
\footnotesize
\begin{itemize}\setlength\itemsep{8pt}
\item Silberschatz, A., Korth, H. F., \& Sudarshan, S. (2014). \textit{Fundamentos de bases de datos} (6.ª ed.). McGraw-Hill.
\item Oracle. MySQL Reference Manual: \href{https://dev.mysql.com/doc/refman/8.0/en/union.html}{\textit{UNION}}, \href{https://dev.mysql.com/doc/refman/8.0/en/intersect.html}{\textit{INTERSECT}} y \href{https://dev.mysql.com/doc/refman/8.0/en/except.html}{\textit{EXCEPT}}.
\item Oracle. MySQL Reference Manual: \href{https://dev.mysql.com/doc/refman/8.0/en/subqueries.html}{\textit{Subqueries}} y \href{https://dev.mysql.com/doc/refman/8.0/en/exists-and-not-exists-subqueries.html}{\textit{Subqueries with EXISTS or NOT EXISTS}}.
\item Oracle. MySQL Reference Manual: \href{https://dev.mysql.com/doc/refman/8.0/en/with.html}{\textit{WITH (Common Table Expressions)}}.
\end{itemize}
\end{frame}
% Cierre UTP
{
  \setbeamertemplate{background canvas}{%
    \includegraphics[width=\paperwidth,height=\paperheight]{gracias-final.pdf}%
  }
  \setbeamertemplate{footline}{}
  \begin{frame}[plain]
  \end{frame}
}

\end{document}
